\documentclass[a4paper,12pt]{article}
\usepackage{fontspec}
\setmainfont{Times New Roman}
\usepackage{polyglossia}
\setdefaultlanguage{russian}
\usepackage[a4paper,left=30mm,right=30mm,top=30mm,bottom=30mm]{geometry}
\usepackage{amsmath}
\usepackage{amssymb}
\usepackage{graphicx}
\usepackage{array}
\usepackage{float}
\usepackage{tikz}
\usepackage{pgfplots}
\usepackage{listings}
\usepackage{setspace}
\usepackage{tocloft}

\pgfplotsset{compat=1.18}
\setstretch{1.15}
\setlength{\parindent}{1.25cm}
\setcounter{secnumdepth}{0}

\renewcommand{\contentsname}{Содержание}
\renewcommand{\cfttoctitlefont}{\hfill\Large\bfseries}
\renewcommand{\cftaftertoctitle}{\hfill\null\par\vspace{1cm}}
\renewcommand{\cftsecleader}{\cftdotfill{\cftdotsep}}
\renewcommand{\cftsubsecleader}{\cftdotfill{\cftdotsep}}
\setlength{\cftbeforesecskip}{0.6em}
\setlength{\cftbeforesubsecskip}{0.2em}

\lstset{
    basicstyle=\ttfamily\small,
    breaklines=true,
    frame=single
}

\begin{document}
\begin{titlepage}
\thispagestyle{empty}
\begin{center}
{\small
МОСКОВСКИЙ ГОСУДАРСТВЕННЫЙ УНИВЕРСИТЕТ\\
имени М.~В.~ЛОМОНОСОВА\\
ФАКУЛЬТЕТ ВЫЧИСЛИТЕЛЬНОЙ МАТЕМАТИКИ И КИБЕРНЕТИКИ
}

\vspace{5cm}

{\Large ОТЧЕТ ПО ЗАДАНИЮ №6}

\vspace{1.2cm}

\begin{minipage}{0.9\textwidth}
\centering
{\large\bfseries
<<Сборка многомодульных программ.\\
Вычисление корней уравнений и определенных интегралов>>
}
\end{minipage}

\vspace{1.2cm}

{\large\bfseries Вариант 9}
\end{center}

\vfill

\begin{flushright}
Исполнитель: студентка 104 группы\\
Ларичева Д.~А.\\[0.7cm]
Преподаватель:\\
Гуляев Д.~А.
\end{flushright}

\vfill

\begin{center}
Москва, 2026
\end{center}
\end{titlepage}

\tableofcontents
\newpage

\section{Постановка задачи}
Целью работы является разработка многомодульной программы на языках C и NASM32 для вычисления площади плоской фигуры, ограниченной тремя кривыми.

В варианте 9 заданы функции:
\[
f_1(x)=\frac{3}{(x-1)^2+1},
\]
\[
f_2(x)=\sqrt{x+0.5},
\]
\[
f_3(x)=e^{-x}.
\]

По индивидуальному варианту необходимо использовать:
\begin{itemize}
    \item точности вычислений: $\varepsilon=10^{-3},10^{-5},10^{-7}$;
    \item тип вычислений: \texttt{float};
    \item методы поиска корней: метод деления пополам и метод Ньютона;
    \item метод вычисления интегралов: метод прямоугольников.
\end{itemize}

Требуется найти точки пересечения кривых, представить площадь фигуры как сумму определенных интегралов и вычислить ее с заданной точностью.

Функции $f_1$, $f_2$, $f_3$ и их производные реализуются на языке NASM32. Остальная часть программы реализуется на языке C. Сборка программы выполняется с помощью \texttt{Makefile}.

\newpage
\section{Математическое обоснование}
Точки пересечения двух кривых находятся как корни уравнения:
\[
f_i(x)=f_j(x).
\]

Для применения численных методов это уравнение приводится к виду:
\[
F(x)=f_i(x)-f_j(x)=0.
\]

В работе используются три уравнения:
\[
f_1(x)-f_3(x)=0,
\]
\[
f_2(x)-f_3(x)=0,
\]
\[
f_1(x)-f_2(x)=0.
\]

Для поиска корней были выбраны следующие отрезки:
\begin{itemize}
    \item $f_1=f_3$: $[-0.5;0]$;
    \item $f_2=f_3$: $[0;0.5]$;
    \item $f_1=f_2$: $[1.5;2.5]$.
\end{itemize}

Выбор отрезков был сделан по графикам функций и проверке смены знака функции $F(x)$ на концах отрезка.

Метод деления пополам основан на следующем свойстве: если функция непрерывна на отрезке $[a,b]$ и выполняется условие
\[
F(a)F(b)<0,
\]
то на этом отрезке существует корень. На каждом шаге вычисляется середина отрезка:
\[
c=\frac{a+b}{2}.
\]

Затем выбирается та половина отрезка, на которой функция меняет знак.

Метод Ньютона задается формулой:
\[
x_{n+1}=x_n-\frac{F(x_n)}{F'(x_n)}.
\]

Для этого нужны производные функций:
\[
f_1'(x)=\frac{-6(x-1)}{((x-1)^2+1)^2},
\]
\[
f_2'(x)=\frac{1}{2\sqrt{x+0.5}},
\]
\[
f_3'(x)=-e^{-x}.
\]

Метод Ньютона при хорошем начальном приближении сходится быстрее метода деления пополам \cite{math1}.

Для вычисления интегралов используется метод прямоугольников:
\[
\int_a^b f(x)\,dx \approx h\sum_{k=0}^{n-1} f(x_k),
\]
где
\[
h=\frac{b-a}{n}.
\]

Количество разбиений увеличивается до тех пор, пока разность между двумя последовательными приближениями не станет меньше заданной точности.

\subsection{Графики функций}
\begin{figure}[H]
\centering
\begin{tikzpicture}
\begin{axis}[
    axis lines=middle,
    xlabel={$x$},
    ylabel={$y$},
    xmin=-1,
    xmax=3,
    ymin=0,
    ymax=3.5,
    grid=major,
    width=12cm,
    height=8cm,
    legend cell align=left,
    legend pos=north east
]
\addplot[blue,thick,domain=-1:3,samples=300] {3/((x-1)^2+1)};
\addlegendentry{$f_1(x)=\frac{3}{(x-1)^2+1}$}
\addplot[red,thick,domain=-0.5:3,samples=300] {sqrt(x+0.5)};
\addlegendentry{$f_2(x)=\sqrt{x+0.5}$}
\addplot[green!60!black,thick,domain=-1:3,samples=300] {exp(-x)};
\addlegendentry{$f_3(x)=e^{-x}$}
\end{axis}
\end{tikzpicture}
\caption{Графики функций варианта 9}
\end{figure}

\newpage
\section{Результаты экспериментов}
В результате работы программы были получены точки пересечения, представленные в таблице.

\begin{table}[H]
\centering
\begin{tabular}{|c|c|c|}
\hline
Кривые & $x$ & $y$ \\
\hline
$f_1=f_3$ & -0.2033349 & 1.2254827 \\
\hline
$f_2=f_3$ & 0.1874113 & 0.8291027 \\
\hline
$f_1=f_2$ & 1.9561526 & 1.5672110 \\
\hline
\end{tabular}
\caption{Координаты точек пересечения}
\end{table}

Площадь фигуры была вычислена как сумма интегралов на промежутках между найденными точками пересечения.

На промежутке от $x_{13}$ до $x_{23}$ область ограничена функциями $f_1$ и $f_3$.

На промежутке от $x_{23}$ до $x_{12}$ область ограничена функциями $f_1$ и $f_2$.

Итоговое значение площади:
\[
S \approx 2.3386023.
\]

\subsection{Количество итераций}
Для метода деления пополам были получены следующие значения:
\begin{table}[H]
\centering
\begin{tabular}{|c|c|}
\hline
Уравнение & Количество итераций \\
\hline
$f_1=f_3$ & 15 \\
\hline
$f_2=f_3$ & 15 \\
\hline
$f_1=f_2$ & 16 \\
\hline
\end{tabular}
\caption{Количество итераций метода деления пополам}
\end{table}

Для метода Ньютона были получены следующие значения:
\begin{table}[H]
\centering
\begin{tabular}{|c|c|}
\hline
Уравнение & Количество итераций \\
\hline
$f_1=f_3$ & 3 \\
\hline
$f_2=f_3$ & 3 \\
\hline
$f_1=f_2$ & 3 \\
\hline
\end{tabular}
\caption{Количество итераций метода Ньютона}
\end{table}

Из таблиц видно, что метод Ньютона для выбранных отрезков сошелся быстрее метода деления пополам.

\subsection{Результаты тестирования}
\begin{table}[H]
\centering
\begin{tabular}{|c|c|c|}
\hline
Тест & Получено & Ожидалось \\
\hline
Бисекция $f_1=f_3$ & -0.2033310 & -0.2033349 \\
\hline
Ньютон $f_1=f_3$ & -0.2033349 & -0.2033349 \\
\hline
Бисекция $f_2=f_3$ & 0.1874161 & 0.1874113 \\
\hline
Ньютон $f_2=f_3$ & 0.1874113 & 0.1874113 \\
\hline
Бисекция $f_1=f_2$ & 1.9561539 & 1.9561525 \\
\hline
\end{tabular}
\caption{Тестирование поиска корней}
\end{table}

\begin{table}[H]
\centering
\begin{tabular}{|c|c|c|}
\hline
Интеграл & Получено & Ожидалось \\
\hline
$\int_0^1 f_1(x)\,dx$ & 2.3561982 & 2.3561945 \\
\hline
$\int_0^2 f_2(x)\,dx$ & 2.3995331 & 2.3995291 \\
\hline
$\int_0^1 f_3(x)\,dx$ & 0.6321140 & 0.6321206 \\
\hline
\end{tabular}
\caption{Тестирование вычисления интегралов}
\end{table}

\newpage
\section{Структура программы и спецификация функций}
Программа состоит из следующих файлов:
\begin{verbatim}
lab6/
├── asm/
│   └── functions.asm
├── include/
│   ├── functions.h
│   └── numerical.h
├── src/
│   ├── main.c
│   └── numerical.c
└── Makefile
\end{verbatim}

\subsection{main.c}
Файл \texttt{main.c} является главным модулем программы. В нем выполняется обработка аргументов командной строки, вызов численных методов и вывод результатов.

Поддерживаются следующие режимы:
\begin{itemize}
    \item вывод площади фигуры;
    \item вывод точек пересечения;
    \item вывод количества итераций;
    \item тестирование поиска корней;
    \item тестирование интегралов.
\end{itemize}

\subsection{numerical.c}
Файл \texttt{numerical.c} содержит реализацию численных методов:
\begin{itemize}
    \item \texttt{root\_bisection} --- метод деления пополам;
    \item \texttt{root\_newton} --- метод Ньютона;
    \item \texttt{integral\_rect} --- метод прямоугольников.
\end{itemize}

\subsection{functions.asm}
Файл \texttt{functions.asm} содержит реализацию функций и их производных на языке NASM32:
\begin{itemize}
    \item \texttt{f1};
    \item \texttt{f2};
    \item \texttt{f3};
    \item \texttt{df1};
    \item \texttt{df2};
    \item \texttt{df3}.
\end{itemize}

Функции вызываются из C-кода по соглашению cdecl.

\newpage
\section{Сборка программы (Make-файл)}
Сборка программы выполняется с помощью \texttt{Makefile}. Используются компилятор \texttt{gcc}, ассемблер \texttt{nasm} и ключ \texttt{-m32} для сборки 32-битной программы.

Основные команды:
\begin{lstlisting}
make
make METHOD=NEWTON
make EPS=1e-3
make EPS=1e-5
make EPS=1e-7
make clean
\end{lstlisting}

Команда \texttt{make} собирает программу с методом деления пополам.

Команда \texttt{make METHOD=NEWTON} собирает программу с методом Ньютона.

Команда \texttt{make EPS=...} позволяет задать точность на этапе сборки.

Команда \texttt{make clean} удаляет объектные файлы и исполняемый файл.

Пример команд компиляции:
\begin{lstlisting}
gcc -m32 -Wall -Wextra -O2 -Iinclude
nasm -f elf32
gcc -m32 -o prog ...
\end{lstlisting}

\newpage
\section{Отладка программы, тестирование функций}
Для проверки корректности были реализованы специальные режимы запуска.

Тестирование поиска корней:
\begin{lstlisting}
./prog --test-root
\end{lstlisting}

Пример вывода:
\begin{lstlisting}
bisection f1=f3: x=-0.2033310, expected about -0.2033349
newton f1=f3: x=-0.2033349, expected about -0.2033349
bisection f2=f3: x=0.1874161, expected about 0.1874113
newton f2=f3: x=0.1874113, expected about 0.1874113
bisection f1=f2: x=1.9561539, expected about 1.9561525
\end{lstlisting}

Тестирование интегралов:
\begin{lstlisting}
./prog --test-integral
\end{lstlisting}

Пример вывода:
\begin{lstlisting}
integral f1[0,1] = 2.3561982, expected about 2.3561945
integral f2[0,2] = 2.3995331, expected about 2.3995291
integral f3[0,1] = 0.6321140, expected about 0.6321206
\end{lstlisting}

Также были проверены параметры:
\begin{itemize}
    \item \texttt{--help};
    \item \texttt{--roots};
    \item \texttt{--iterations};
    \item \texttt{--test-root};
    \item \texttt{--test-integral}.
\end{itemize}

\newpage
\section{Программа на Си и на Ассемблере}
Программа реализована на языках C и NASM32.

На языке ассемблера реализованы функции варианта и их производные. На языке C реализованы численные методы и основная логика программы.

Исходные тексты программы находятся в архиве, приложенном к отчету.

В архив входят:
\begin{itemize}
    \item \texttt{main.c};
    \item \texttt{numerical.c};
    \item \texttt{functions.asm};
    \item \texttt{functions.h};
    \item \texttt{numerical.h};
    \item \texttt{Makefile}.
\end{itemize}

\newpage
\section{Анализ допущенных ошибок}
Во время выполнения работы возникали следующие ошибки:
\begin{itemize}
    \item сначала был выбран неправильный набор функций;
    \item были неправильно подобраны интервалы поиска корней;
    \item возникали ошибки при работе с FPU;
    \item возникали ошибки при подключении NASM-модуля к C-программе;
    \item метод Ньютона требовал аккуратного выбора начального приближения.
\end{itemize}

После исправления функций, интервалов и ассемблерной части программа начала корректно вычислять точки пересечения и площадь фигуры.

\newpage
\section{Литература}
\begin{thebibliography}{99}
\bibitem{math1}
Ильин В.~А., Садовничий В.~А., Сендов Бл.~Х.
Математический анализ. Т.~1.
--- М.: Наука, 1985.

\bibitem{method}
Трифонов Н.~П., Пильщиков В.~Н.
Задания практикума на ЭВМ.
\end{thebibliography}

\end{document}